% The L-symbols along a differentiable symmetric path.
% Source: GLM23, arXiv:2203.12563v3, REsubmission.tex, lines 1352--1578
% (MPSs with inequivalent L-symbols belong to different phases).
\section{The multiplicity symbols along a symmetric path}
\label{sec:glm_l_symbol_path}

Garre-Rubio, Lootens and Moln\'ar show that a differentiable path of
symmetric states cannot change the gauge class of the $L$-symbols
\cite[Section~5]{GarreRubio2023MPOSym}. Along the path the operator blocks,
the fusion maps and the bond dimensions are fixed. The derivative of each
action tensor is assumed to lie in the span of the action tensors at the
same point. Under these assumptions the first-order change of the
coefficient matrices is an infinitesimal gauge transformation. This section
proves that statement. The left inverses of the action tensors are not
assumed to vary differentiably.

\begin{lemma}[A first-order quotient rule]
    \leanok
    \label{lem:glm_l_path_quotient_rule}
    \lean{Matrix.hasDerivAt_mul_entries,Matrix.hasDerivAt_trace_of_entries,
        Matrix.differentiableAt_det_of_entries,
        Matrix.differentiableAt_inv_of_entries,
        Matrix.hasDerivAt_entries_of_eventually_eq_mul}
    Let $R(\gamma)$, $L(\gamma)$ and $P(\gamma)$ be finite complex matrices
    depending on a real parameter, with $P$ square. Suppose $R=LP$ on a
    neighbourhood of $\gamma_0$ and $P(\gamma_0)=1$. Suppose also that the
    entries of $R$ and $P$ are differentiable at $\gamma_0$ with derivatives
    $R'$ and $P'$. Then the entries of $L$ are differentiable at $\gamma_0$,
    with derivative $R'-L(\gamma_0)P'$.
\end{lemma}
\begin{proof}
    \leanok
    The determinant of $P$ is a polynomial in its entries. It is therefore
    differentiable at $\gamma_0$, and it is nonzero near $\gamma_0$ because
    $\det P(\gamma_0)=1$. By Cramer's rule the entries of $P^{-1}$ are
    differentiable at $\gamma_0$. Near $\gamma_0$ we have $L=RP^{-1}$, so
    $L$ is differentiable there. The product rule applied to $R=LP$ gives
    $R'=L'P(\gamma_0)+L(\gamma_0)P'=L'+L(\gamma_0)P'$.
\end{proof}

\begin{theorem}[The change of the $L$-symbols along a path is a gauge
        transformation]
    \leanok
    \label{thm:glm_l_path_invariance}
    \lean{MPOTensor.hasDerivAt_actionLMatrix}
    \uses{def:glm_boundary_action_l_matrix}
    Fix operator blocks $O_a$ of positive bond dimensions with exact
    biorthogonal fusion maps. Let $A_x(\gamma)$ be state blocks of fixed
    positive bond dimensions $D_x$. Assume that for $\gamma$ near $\gamma_0$
    they are individually injective and pairwise inequivalent under a
    virtual gauge and a nonzero scalar. Assume also that they carry exact
    biorthogonal action maps $V_{ax}^{y,i}(\gamma)$ and
    $\widehat V_{ax}^{y,i}(\gamma)$. Suppose every entry of each analysis
    map is differentiable at $\gamma_0$, with
    \begin{align}
        \frac{d}{d\gamma}V_{ax}^{y,i}(\gamma)\Big|_{\gamma_0}
         & =\sum_{i'}(X_{ax}^y)_{ii'}V_{ax}^{y,i'}(\gamma_0)
        \label{eq:glm_l_path_span}
    \end{align}
    for some square matrices $X_{ax}^y$. Then every entry of the
    coefficient matrices $L_{abx}^y(\gamma)$ is differentiable at
    $\gamma_0$. Writing $L$ for $L_{abx}^y(\gamma_0)$, the derivative is
    \begin{align}
        \frac{d}{d\gamma}L_{abx}^y(\gamma)[(z,i,j),(c,k,\mu)]\Big|_{\gamma_0}
         & =\sum_{i'}(X_{az}^y)_{ii'}L[(z,i',j),(c,k,\mu)]
        +\sum_{j'}(X_{bx}^z)_{jj'}L[(z,i,j'),(c,k,\mu)]
        \notag                                                   \\
         & \quad-\sum_{k'}L[(z,i,j),(c,k',\mu)](X_{cx}^y)_{k'k}.
        \label{eq:glm_l_path_derivative}
    \end{align}
    This is the invariance step of \cite[Section~5]{GarreRubio2023MPOSym},
    whose text reads, with the pictured expansion of the action tensors
    written as the bracketed formula:
    \begin{quote}
        Since one can always find a left inverse tensor such that it is
        continuous we can find without loss of generality a basis change,
        where $V_{a,x}^{y,i}$ is also continuous (using the fact that the
        action tensors are linearly independent). Let us absorb this basis
        change and assume (see below for generalization) that the action
        tensors are differentiable:
        [$V^{i}(\gamma+d\gamma)=V^{i}(\gamma)
            +\sum_{i'}(X_{ax}^y)_{i'}^{i}V^{i'}(\gamma)\,d\gamma$], where we
        have spanned the term in $d\gamma$ using the orthogonal set of action
        tensors $\{V_{a,x}^{y,i},i=1,\cdots,M_{a,x}^y\}$ of the subspace of
        the block $y$. [\dots] Therefore
        \[
            ({L'}_{abx}^y)_{c\mu k}^{zij}=\sum_{\hat i,\hat j,\hat k}
            (L_{abx}^y)^{z\hat i\hat j}_{c\mu\hat k}\big[
                \delta_{i,\hat i}\delta_{j,\hat j}\delta_{k,\hat k}
                +(X_{bx}^z)_{\hat j}^{j}\delta_{i,\hat i}\delta_{k,\hat k}
                +(X_{az}^y)_{\hat i}^{i}\delta_{j,\hat j}\delta_{k,\hat k}
                -(X_{cx}^y)^{\hat k}_{k}\delta_{i,\hat i}\delta_{j,\hat j}\big],
        \]
        which means that the ${L}'$-symbol is in the same class as the
        ${L}$-symbol since they are related by a (differentiable) gauge
        transformation.
    \end{quote}
    The source's $(X_{ax}^y)_{i'}^{i}$ is the entry $(X_{ax}^y)_{ii'}$ here.
    The hypotheses are those of the source: the symmetry operators and
    fusion maps are the same along the path, the blocks are injective, the
    bond dimension does not change, and (\ref{eq:glm_l_path_span}) holds.
    The source's assertion that the new $L$-symbols are differentiable is
    proved here, not assumed.
\end{theorem}
\begin{proof}
    \leanok
    \uses{thm:glm_boundary_action_l_exists,lem:glm_l_path_quotient_rule}
    Fix $a,b,x,y$. Let $S_q$ be the fusion-then-action synthesis maps at
    $\gamma_0$, indexed by $q=(c,k,\mu)$. Let $H_{{\rm seq},p}(\gamma)$ and
    $H_{{\rm fus},q}(\gamma)$ be the sequential and fusion-then-action
    analysis maps at $\gamma$, with $p=(z,i,j)$. The defining relation
    (\ref{eq:glm_boundary_action_l_defining}) at $\gamma$, restricted to the
    final block $y$, reads
    $H_{{\rm seq},p}(\gamma)=\sum_qL(\gamma)_{pq}H_{{\rm fus},q}(\gamma)$.
    Multiplying on the right by $S_{q'}$ and taking normalized traces
    gives $R(\gamma)=L(\gamma)P(\gamma)$, where
    $R(\gamma)_{pq'}=D_y^{-1}\operatorname{tr}(H_{{\rm seq},p}(\gamma)S_{q'})$
    and
    $P(\gamma)_{qq'}=D_y^{-1}\operatorname{tr}(H_{{\rm fus},q}(\gamma)S_{q'})$.
    Biorthogonality at $\gamma_0$ gives
    $H_{{\rm fus},q}(\gamma_0)S_{q'}=\delta_{qq'}1$, so $P(\gamma_0)=1$.
    The maps $H_{\rm fus}$ are linear and $H_{\rm seq}$ bilinear in the
    action analysis maps. By (\ref{eq:glm_l_path_span}),
    \begin{align}
        R'_{pq'}         & =\sum_{i'}(X_{az}^y)_{ii'}L_{(z,i',j)q'}
        +\sum_{j'}(X_{bx}^z)_{jj'}L_{(z,i,j')q'},
        \notag                                                              \\
        P'_{(c,k,\mu)q'} & =\sum_{k'}(X_{cx}^y)_{kk'}\delta_{(c,k',\mu)q'}.
        \notag
    \end{align}
    Lemma~\ref{lem:glm_l_path_quotient_rule} gives the derivative
    $R'-LP'$, which is (\ref{eq:glm_l_path_derivative}).
\end{proof}

\begin{corollary}[The derivative is tangent to a gauge orbit]
    \leanok
    \label{cor:glm_l_path_gauge_tangent}
    \lean{MPOTensor.hasDerivAt_actionLMatrix_gaugeCurve}
    \uses{thm:glm_l_path_invariance,lem:glm_mixed_gauge_coefficients}
    Apply the gauge law (\ref{eq:glm_mixed_gauge_l}) to the coefficients
    $L=L_{abx}^y(\gamma_0)$ with $X_{ax}^y(t)=1+tX_{ax}^y$, with
    $(\widehat X_{cx}^y(t))_{kk'}=(1-tX_{cx}^y)_{k'k}$, the transposed
    inverse to first order, and with trivial fusion gauge. The resulting
    curve of coefficients has velocity at $t=0$ equal to the right side of
    (\ref{eq:glm_l_path_derivative}). Thus the first-order change of the
    $L$-symbols along the path is the infinitesimal gauge transformation
    generated by the matrices $X_{ax}^y$, and the class of the $L$-symbols
    is unchanged to first order.
\end{corollary}
\begin{proof}
    \leanok
    The curve is a polynomial in $t$. Its derivative at $t=0$ collects the
    three terms linear in $t$, one for each gauge factor.
\end{proof}
